\label{sec:introduction}

\textit{Thornburg v.\ Gingles}~\citep{gingles1986} established three
preconditions---commonly called ``the Gingles prongs''---that a minority
group must satisfy to state a vote-dilution claim under Section 2 of the
Voting Rights Act~\citep{vra1965}:

\begin{enumerate}
  \item \textbf{Prong 1 (Geographic Compactness):} The minority group must be
    ``sufficiently large and geographically compact to constitute a majority
    in a single-member district.''
  \item \textbf{Prong 2 (Political Cohesion):} The minority group must be
    ``politically cohesive.''
  \item \textbf{Prong 3 (Majority Bloc Voting):} The majority must ``vote
    sufficiently as a bloc to enable it \ldots\ usually to defeat the
    minority's preferred candidate.''
\end{enumerate}

Together, the three prongs establish that a minority group can form a
majority district (Prong 1), would use that district to elect its preferred
candidate (Prong 2), and is currently prevented from doing so because the
majority votes as a bloc against minority-preferred candidates (Prong 3).
Satisfying all three prongs does not guarantee a Section 2 violation; a
``totality of circumstances'' inquiry follows, but in practice the Gingles
prongs are the dispositive hurdle.

Litigation under Section 2 has evolved substantially since \textit{Gingles}.
\textit{Allen v.\ Milligan}~\citep{allen2023} reaffirmed the core Gingles
framework and rejected Alabama's invitation to replace it with a race-neutral
comparator methodology.  \textit{Louisiana v.\ Callais}~\citep{callais2026}
added a specific methodological requirement for Prong 3: expert witnesses must
\emph{disentangle} racial bloc voting from partisan bloc voting using a
regression with a partisan control variable, because post-\textit{Rucho}
states have increasingly defended against VRA claims by arguing that observed
racial bloc voting is actually partisan and therefore constitutionally
permissible~\citep{rucho2019}.

This paper presents a systematic quantitative methodology for satisfying all
three Gingles prongs using the \texttt{redist} algorithmic redistricting
system.  The methodology is designed for use in expert witness contexts.
The contributions are:

\begin{enumerate}
  \item A principled operationalization of Prong 1 compactness using
    VRASection's geographic alignment score (Section~2).
  \item A regression-based Prong 2 analysis using \texttt{redist analyze
    --types bloc-voting} with primary election data (Section~3).
  \item A post-\textit{Callais} Prong 3 methodology using WLS+HC3+Holm
    regression with a partisan control variable, with confidence intervals
    from bootstrap resampling (Section~4).
  \item A worked example using the Alabama AL-07 district from
    \textit{Allen v.\ Milligan} (Section~5).
  \item An expert witness checklist mapping \texttt{redist} outputs to each
    Gingles prong (Section~7).
\end{enumerate}

The methodology is algorithm-agnostic in its Prong 2 and Prong 3
components---the WLS+HC3+Holm regression can be applied to any district
configuration, not just algorithmic ones.  The distinctive contribution of
the \texttt{redist} system is Prong 1: the VRASection alignment score provides
a principled, reproducible, and legally defensible threshold for geographic
compactness that replaces the ad-hoc visual inspection that has historically
dominated expert practice.
